\section{Introduction}

This chapter sets out the broader research context in which the present thesis is situated. It examines the overarching motivations that drive the work as a whole, before detailing the specific aims and rationale for each of the subsequent chapters. While each chapter addresses a different facet of the overall research problem, they are unified by a central concern: the conservation of musical instruments using digital methods, and the parallel requirement to conserve the very digital tools on which such conservation efforts rely.

The research context naturally divides into two complementary domains:

\begin{enumerate}
    \item \textbf{Digital musical instrument conservation}: the use of measurement, modelling, and interaction design to document, preserve, and—where possible—replicate the acoustical and tactile properties of musical instruments.
    \item \textbf{Digital conservation more generally}: the safeguarding of the digital artefacts—data, code, formats, and workflows—used to perform the conservation itself, ensuring their long-term accessibility and reproducibility.
\end{enumerate}

Musical instrument conservation addresses the specific challenges inherent in preserving objects that are not merely visual artefacts, but functional tools whose significance lies equally in their physical form, acoustical output, and the embodied interaction they enable. Digital conservation, in turn, deals with the processes, protocols, and infrastructures required to ensure that the software, datasets, and digital surrogates created in support of cultural heritage remain available, interpretable, and usable far into the future.

The combination of these two domains gives rise to a layered conservation problem. At one level, we aim to capture the sonic, structural, and functional essence of an instrument through non-destructive measurement and modelling. At another, we must acknowledge that the tools used in this process—measurement software, simulation codes, repositories—are themselves ephemeral and require active management to prevent their loss.

In addressing this second point, the work aligns itself with the principles of the Open Science movement, which since around 2016 has increasingly shaped research practice across disciplines. Central to this movement are commitments to transparency, reproducibility, and sustainability, exemplified by frameworks such as the FAIR principles. Although originally conceived for scientific data, these principles can be applied with equal force to the digital products of cultural heritage work.

This chapter therefore proceeds in two major sections. First, it examines the context of musical instrument conservation, identifying the principal challenges and limitations—both technical and institutional—that shape current practice. Second, it defines the term \emph{digital conservation} as used in this thesis, discussing the challenges of integrating digital tools into cultural heritage workflows and outlining a methodology—grounded in FAIR principles—for doing so in a sustainable, reproducible way.

\section{Musical Instrument Conservation}

Musical instruments present a unique conservation challenge. Unlike most museum artefacts, they are \emph{functional} objects, whose full significance is not exhausted by their material appearance. Their value lies not only in their historical and aesthetic qualities, but also in their acoustical behaviour and tactile response under the hands of a player. Preserving this dual identity poses a constant tension between two imperatives:

\begin{enumerate}
    \item \textbf{Preservation of material integrity}: avoiding physical deterioration, often by restricting handling, controlling environmental conditions, and minimising mechanical stress.
    \item \textbf{Preservation of functional identity}: maintaining the instrument in playable condition, or at least preserving the information necessary to reproduce its acoustical and mechanical behaviour.
\end{enumerate}

These imperatives often conflict. For example, stringed keyboard instruments—harpsichords, clavichords, fortepianos—rely on delicate wooden structures and organic materials that inevitably degrade over time. Keeping such instruments in playing condition can accelerate wear, while restricting their use may protect the object but at the cost of losing an understanding of its sonic and mechanical properties.

Traditionally, museums have taken one of two approaches. In some cases, instruments are maintained for performance, accepting controlled wear as a trade-off for keeping their living function alive. In others, they are retired from use, displayed as static exhibits, and interpreted through recordings or replicas. Both approaches involve compromises, and neither fully addresses the challenge of preserving both the tangible and intangible aspects of an instrument.

Digital technology offers a potential third path. Advances in high-resolution geometry capture, acoustical measurement techniques, and physical modelling enable the creation of detailed digital surrogates that document an instrument’s structure, materials, and sonic response. These surrogates can be used to drive real-time simulations or digitally controlled replicas, allowing users to interact with a convincing representation of the instrument without risking the original.

Yet such digital surrogates are only as robust as the measurement, modelling, and preservation processes that underpin them. They depend on accurate data acquisition, transparent processing methods, and sustainable storage solutions—connecting the task of musical instrument conservation directly to the broader field of digital conservation.

\section{Digital Conservation and FAIR Principles}

In the scope of this thesis, \emph{digital conservation} refers to the deliberate and systematic preservation of digital assets so that they remain accessible, interpretable, and usable over extended time scales. Here, “digital assets” encompass any form of digital material created in the course of conservation work: raw measurement data, processed datasets, source code, compiled binaries, CAD designs, 3D models, and associated documentation.

Long-term survival of such assets depends on more than just storing files. File formats may become obsolete, hardware and operating systems change, and software dependencies can cease to be available. Without clear metadata, persistent identifiers, and thorough documentation, even well-preserved files may become effectively useless.

The FAIR principles, articulated by \textcite{wilkinson_fair_2016}, provide a concise and actionable framework for addressing these challenges. Their four components are:

\begin{itemize}
    \item \textbf{Findability}: Digital assets should be described with rich, standardised metadata and assigned globally unique, persistent identifiers (such as DOIs) so they can be reliably located by humans and machines.
    \item \textbf{Accessibility}: Assets should be retrievable using open, universally implementable communication protocols, with clear access conditions.
    \item \textbf{Interoperability}: Assets should use formal, accessible, and broadly applicable languages, formats, and vocabularies, so that they can be integrated with other data or tools with minimal effort.
    \item \textbf{Reusability}: Assets should be released under explicit, accessible usage licences and accompanied by detailed documentation, enabling meaningful reuse in new contexts.
\end{itemize}

While originally aimed at research data management, these principles are equally applicable to the products of cultural heritage conservation, where digital outputs may have a life far beyond their initial project context. Applying FAIR in this domain means ensuring that measurement datasets, simulation codes, and digital models are stored in formats that will remain usable; that their provenance and dependencies are recorded; and that they are shared in ways that encourage and enable reuse.

Practical adherence to FAIR in this thesis has included: placing all digital work under version control using \texttt{git}; hosting repositories publicly on platforms such as GitHub to maximise accessibility; archiving formal “releases” on Zenodo with assigned DOIs; and licensing all software under GPLv3 to guarantee the freedoms to run, study, share, and modify, while ensuring that derivative works remain open.

A “release” is here defined as a self-contained, functional version of an asset that has been used in the course of the research and is fixed at a particular point in time. Digital assets evolve—features are added, bugs are fixed, formats are refined—but releases provide immutable reference points that can be cited, reproduced, and compared in the future.

\subsection{Methods for Reproducibility}

Reproducibility is a cornerstone of both scientific research and sustainable digital conservation. As noted by \textcite{perkel_challenge_2020}, while perfect computational reproducibility over time is rarely achievable, deliberate design choices can maximise the likelihood that others (including one’s future self) can replicate results. This entails:

\begin{itemize}
    \item Encoding all data processing and analysis in code rather than point-and-click interfaces.
    \item Thoroughly documenting code structure, inputs, outputs, and dependencies.
    \item Recording key parameters, such as random seeds, that influence results.
    \item Implementing test suites to validate functionality.
    \item Providing master scripts to automate complete workflows.
    \item Archiving stable releases in long-term repositories with persistent identifiers.
    \item Using containerisation or environment managers to capture execution environments.
\end{itemize}

These practices are embedded in the methodology of this thesis, with each major software output subjected to version control, documented in detail, and archived in a citable form.

\subsection{Addressing Open Science, FAIR, and Software Sustainability}

In light of the above, the outputs of this thesis have been developed with explicit attention to versioning, accessibility, licensing, and archiving. Every codebase has been managed in a \texttt{git} repository, hosted publicly to encourage engagement, and released under GPLv3 to ensure that any modifications remain equally open. All formal releases have been archived on Zenodo with assigned DOIs, providing a fixed, citable reference point.

This approach aligns with recommendations from the Software Sustainability Institute \cite{hong_how_2019} and is intended to ensure that the thesis outputs remain accessible and functional for as long as possible, regardless of changes in the technological landscape.

\subsection{Aurora: A Case Study in Software Sustainability}

During the course of this PhD, the passing of Professor Angelo Farina brought renewed attention to the fragility of even widely used research software. Farina, a prolific academic in the field of acoustics, was the author of numerous software tools that have become foundational in acoustic measurement and analysis. Among these, \emph{Aurora} stands as a prominent example—not only as a software package of enduring technical value, but also as a tool that has been integral to digital cultural heritage projects for decades \cite{fausti_acoustic_2000,frey_acoustical_2013,almagropastor_assessing_2021}.

Originally created in 1995 \cite{farina_auralization_1995} as an extension to Cool Edit \cite{tronchin_acoustics_1997}—a commercial audio editor that would later evolve into Adobe Audition \cite{farina_listening_2004}—Aurora was initially distributed alongside Farina’s \emph{Ramsete} software for 3D building acoustics simulation \cite{farina_aurora_1995}. Distributed as shareware from \url{http://aurora.ramsete.com/Aurora} and later from \url{http://aurora-plugins.com/}, Aurora simplified the process of sine sweep room impulse response measurement \cite{farina_aurora_2000} at a time when such workflows were laborious and error-prone.

Over the next decade, Aurora underwent continuous development, with successive refinements improving usability, measurement accuracy, and integration with other acoustic analysis tools. In a significant shift, the software was later rewritten by Simone Campanini for the free and open source audio editor Audacity \cite{farina_new_2009}, marking the first time Aurora could be run entirely within an open environment. This transition ended the “Russian Roulette” shareware model described by Farina \cite{farina_registration_2007} and opened the way for broader community engagement.

Despite its longevity and widespread citation—Aurora appears in research as recent as 2024 \cite{fontana_highly_2024,popp_speech_2024,denti_pan_2024}—its sustainability has been undermined by the absence of a formal, open-source release history. While the Audacity port’s source code was made available, the original Adobe Audition version remains closed. Without formal software citation practices, it is often impossible to determine which Aurora version was used in a given study, making it difficult to reproduce results or assess the impact of code changes on past research outcomes.

This recognition led to a collaborative preservation effort involving Campanini and Farina. The Audacity variant of Aurora was placed in a public \texttt{git} repository, licensed under GPLv3, and archived with a DOI \cite{farina_aurora_2025}, aligning it with the principles advocated by the Software Sustainability Institute \cite{hong_how_2019}. While this does not recover the full development history, it provides a fixed, citable reference and ensures that the code remains available for future work.

Aurora’s technical integration into Audacity via its experimental “module” system—an under-documented feature requiring co-compilation with Audacity’s own source—offers both benefits and risks. The benefit is straightforward: a user can download a single, functioning audio editor environment with Aurora pre-integrated. The risk is equally clear: if Audacity’s source code changes substantially, the Aurora module may no longer compile or function, leaving the software inoperable.

The acquisition of Audacity by Muse Group in 2021 \cite{prokoudine_ultimate_2021} increases the likelihood of such breaking changes, given the pace of ongoing development. Without active maintenance and documentation, Aurora risks slipping into obsolescence, its binaries inaccessible and its functionality preserved only in scattered publications and personal archives.

This predicament is not unique. Similar vulnerabilities can be seen in the case of \emph{YIN}, a widely cited pitch detection algorithm by Alain de Cheveigné \cite{cheveigne_yin_2002}. Originally available from IRCAM at \url{http://www.ircam.fr/pcm/cheveign/sw/yin.zip}, YIN is now hosted informally at \url{http://audition.ens.fr/adc/sw/yin.zip} \cite{alain_cancellation_2023}, but still references the defunct IRCAM link in its source code. Without institutional archiving and DOI assignment, such software remains one personnel change or server migration away from disappearance.

Aurora and YIN illustrate a critical point: neither the popularity of a tool nor its status as “free” or “open source” guarantees its survival. Preservation demands active institutional and community support, explicit licensing, and formal archival processes.

\subsection{Software Citation}

Software citation is a key enabler of reproducibility and credit in research. Proper citation practices identify the exact software and version used, the persistent location where it can be accessed, and the licence governing its use. In the absence of such practices, even widely used tools may leave behind a fragmented and unreliable record of their role in research.

In this thesis, every software output is accompanied by a formal citation, including a DOI where applicable, and is documented with sufficient metadata to satisfy FAIR principles. This ensures that future researchers can both locate the exact version used and evaluate its relevance to their own work.

\section{Definitions}

\begin{itemize}
    \item \textbf{CAD}: Computer-Aided Design. In this context, CAD files are used for the design of printed circuit boards (PCBs) and three-dimensional printed components.
    \item \textbf{PCB}: Printed Circuit Board. A physical substrate—typically a non-conductive plate with copper traces—used to mechanically support and electrically connect electronic components.
    \item \textbf{FAIR}: The four foundational principles of Findability, Accessibility, Interoperability, and Reusability for digital assets \cite{wilkinson_fair_2016}.
    \item \textbf{DOI}: Digital Object Identifier. A persistent, globally unique identifier used to locate and cite digital resources.
    \item \textbf{Digital Asset}: Any form of digital data created and used in the course of research, including but not limited to source code, binaries, CAD files, images, datasets, and documentation.
    \item \textbf{FOSS}: Free and Open Source Software, as defined by the GNU Project’s Four Essential Freedoms \cite{gnu_philosophy_2025}, granting users the rights to run, study, share, and modify software, with source code publicly available.
\end{itemize}

\subsection{Breakdown of the Thesis}

The thesis addresses two distinct stages of the conservation pipeline. \textbf{Chapter~\ref{chapter:keyboard}} addresses the post-digitisation phase, asking: how can interaction with a digitised musical instrument be meaningfully facilitated? \textbf{Chapter~\ref{chapter:magpie}} addresses the pre-digitisation phase, focusing on material analysis of a historical instrument to inform conservation decisions \cite{barclay_care_1997}.

\subsubsection{Chapter 2: Interface}

This chapter documents the design, prototyping, and fabrication of a replica harpsichord keyboard with an early Italian layout, augmented with optical electronic sensors for interactive use with digital instruments. The work responds to the challenge of preserving the tactile and functional aspects of historical instruments in museum contexts, where direct interaction with originals is often prohibited.

By leveraging technologies such as 3D printing, microcontroller-based sensing, and real-time digital synthesis, the project aims to balance historical authenticity with public engagement. The resulting keyboard was deployed at Museo San Colombano, Bologna, as part of an exhibition enabling visitors to interact with digitised instruments in the collection. This deployment forms a case study in how digital interfaces can extend the experiential authenticity of cultural heritage \cite{trant_when_1999, pine_museums_2007} while safeguarding fragile artefacts.

\subsubsection{Chapter 3: MAGPIE}

This chapter presents MAGPIE: an open-source software tool for plate vibration analysis, material property estimation, and acoustics education. The software incorporates numerical models, parameter estimation routines, and visualisation tools, designed for reproducibility and extensibility. The project demonstrates how FOSS tools can lower barriers to advanced analysis techniques in both research and educational contexts.

\subsubsection{Chapter 4: Applications}

Chapter 4 integrates the work of Chapters 2 and 3 by embedding MAGPIE-driven numerical models into the augmented keyboard interface. Sensor data from the keyboard are used to drive bespoke synthesis algorithms informed by measurements of historical harpsichords, incorporating specific mechanical details such as jack staggering and register-dependent plucking dynamics. The chapter also documents external adoption of the tools, including use by the NEMUS project and integration into custom interactive installations.

\subsubsection{Chapter 5: Future}

The concluding chapter considers future directions for both digital musical instrument conservation and the preservation of conservation tools. It reflects on the evolving role of open-source methods in cultural heritage and identifies emerging challenges such as maintaining long-term engagement with archived codebases and developing cross-institutional standards for digital asset management.

\section{Questions Addressed}

The research seeks to answer the following questions:

\begin{itemize}
    \item How can the physical experience of playing a historical musical instrument be preserved in a museum context while protecting the original object?
    \item In what ways can the visual and tactile elements of an instrument’s interface be leveraged to support this preservation?
    \item How can digital technologies, particularly FOSS tools, be applied to support the analysis, conservation, and presentation of musical instruments?
\end{itemize}

\section{Goals and Methodology}

The work adopts an exploratory approach to identify and address three key areas where FOSS digital tools can enhance the conservation of musical instruments:

\begin{enumerate}
    \item \textbf{Analysis}: Processing and interpretation of measurement data to understand material properties without damaging the original artefact.
    \item \textbf{Interaction}: Designing interfaces for digitised instruments that preserve the tactile and performative aspects of the original.
    \item \textbf{Application}: Integrating and deploying these tools in real-world conservation and exhibition contexts, assessing their effectiveness and identifying opportunities for further development.
\end{enumerate}

Each area is addressed through a combination of literature review, technical development, empirical testing, and public deployment, with all outputs documented, versioned, and archived in accordance with FAIR principles.


